%%%%%%%%%%%%%%%%%%%%%%%%%%%%%%%%%%%%%%%%%%%%%%%%%%%%%%%%%%%%%%%%%%%%%%%%%%%%%%%%%%
% Session 19 opening (Sep 2026): the Netflix Prize story as a case for "a model is only useful once it is deployed".
% Session 19 only; ml_predictive_analytics.tex keeps its own two-frame version of the story for the Deployment seminar.
\begin{frame}[fragile]\frametitle{}
\begin{center}
{\Large A Cautionary Tale: The Netflix Prize}

{\tiny (Ref: Netflix Technology Blog, Netflix Recommendations: Beyond the 5 Stars, 2012)}
\end{center}
\end{frame}

%%%%%%%%%%%%%%%%%%%%%%%%%%%%%%%%%%%%%%%%%%%%%%%%%%%
\begin{frame}[fragile]\frametitle{Why Netflix Cares About Predictions}
\begin{itemize}
\item Netflix earns money when members keep watching: a good recommendation keeps them subscribed
\item Recommendations drive roughly three-quarters of what its members watch
\item In 2006 its recommender (Cinematch) predicted the star rating a member would give a film
\item A better rating prediction means a better recommendation
\end{itemize}
\end{frame}

%%%%%%%%%%%%%%%%%%%%%%%%%%%%%%%%%%%%%%%%%%%%%%%%%%%
\begin{frame}[fragile]\frametitle{The \$1 Million Contest}
\begin{itemize}
\item October 2006: Netflix released about 100 million ratings and opened a public contest
\item The goal: beat Cinematch's prediction error by 10\%
\item The prize: \$1 million
\item Thousands of teams from around the world entered
\end{itemize}
\end{frame}

%%%%%%%%%%%%%%%%%%%%%%%%%%%%%%%%%%%%%%%%%%%%%%%%%%%
\begin{frame}[fragile]\frametitle{The Winning Model}
\begin{itemize}
\item It took almost three years to cross the 10\% mark
\item No single model got there: teams merged, and blended their models into ensembles
\item In 2009 the team BellKor's Pragmatic Chaos won
\begin{itemize}
\item A rival ensemble scored the same, but submitted 20 minutes later
\end{itemize}
\item The winning solution blended hundreds of predictive models
\end{itemize}
\end{frame}

%%%%%%%%%%%%%%%%%%%%%%%%%%%%%%%%%%%%%%%%%%%%%%%%%%%
\begin{frame}[fragile]\frametitle{The Model That Never Shipped}
\begin{itemize}
\item Netflix put two methods from the early contest years into production: SVD (matrix factorization) and a Restricted Boltzmann Machine
\item It never ran the winning ensemble in production
\item In Netflix's own words, the extra accuracy ``did not seem to justify the engineering effort needed to bring them into a production environment''
\item The world had also moved on: Netflix was shifting from DVD-by-mail star ratings to streaming, and the data changed
\end{itemize}

\begin{block}{Intuition}
Winning a leaderboard and running a service are different jobs. A tiny gain in accuracy is worth little
if the model is too complex to run, retrain and watch, or if the data it learned from no longer matches reality.
\end{block}
\end{frame}

%%%%%%%%%%%%%%%%%%%%%%%%%%%%%%%%%%%%%%%%%%%%%%%%%%%
\begin{frame}[fragile]\frametitle{Lessons for This Session}
\begin{itemize}
\item Accuracy is not the finish line: a model creates value only when it is deployed and runs reliably
\item Complexity has a running cost: every extra model must be served, updated and monitored
\item Data changes, so a deployed model must be watched (drift, later in this session)
\item Governance matters: a planned second contest was cancelled in 2010 after privacy complaints about the released data
\end{itemize}
Next: what it takes to put a model into production.
\end{frame}

%%%%%%%%%%%%%%%%%%%%%%%%%%%%%%%%%%%%%%%%%%%%%%%%%%%
\begin{frame}[fragile]\frametitle{Quick Check: The Netflix Prize}
\begin{block}{Think About It}
The winning team beat Netflix's 10\% target, yet Netflix did not run the winning ensemble in production. What is the best explanation?
\begin{enumerate}
  \item[A.] The extra accuracy was small compared with the engineering effort to run it, and the business and data had moved on
  \item[B.] The ensemble turned out to be less accurate than Cinematch when tested on live members
  \item[C.] Netflix was not allowed to use the winning models after the contest
  \item[D.] Ensembles cannot be deployed, only single models can
\end{enumerate}
\end{block}
\end{frame}

%%%%%%%%%%%%%%%%%%%%%%%%%%%%%%%%%%%%%%%%%%%%%%%%%%%
\begin{frame}[fragile]\frametitle{Quick Check: The Netflix Prize}
\begin{block}{Answer}
\textbf{Correct: A.} Netflix said the accuracy gains it measured did not seem to justify the engineering effort needed to bring the ensemble into production, and its business was shifting from DVD ratings to streaming. Lower accuracy (B) and legal limits (C) are not the reasons Netflix gave, and ensembles can be deployed (D): the cost is running, retraining and monitoring many models.
\end{block}
\end{frame}